% GENERATED FILE. Edit canonical lessons, then run npm run generate:books.
% canonical-source-hash: fnv1a64:aba0b6e610b7be15
% canonical-lessons: PT-C128-pesado, PT-C128-magro, PT-C128-contente, PT-C128-fim-de-semana, PT-C128-data

\chapter{Heavy, Thin, Pleased, Weekend, Date}
\label{ch:pt-heavy-thin-pleased-weekend-date}
\hlchaptermodality{128}

\begin{chapteropening}
\textbf{By the end of this chapter:} \emph{I can say heavy, thin, pleased, the weekend and a date.}

Complete the last lesson of chapter 128: I can say heavy, thin, pleased, the weekend and a date.
\end{chapteropening}

\section[pesado]{pesado — heavy}

\label{lesson:PT-C128-pesado}



\begin{quote}
Before the new one: say the Portuguese for salty, then the Portuguese for light.
\end{quote}

\subsection*{You'll want to know: pesado}
\textbf{pesado} — \textquotedblleft{}heavy\textquotedblright{}.

\begin{cousinweb}
\textbf{pesado} is the participle of the verb \textquotedblleft{}to weigh\textquotedblright{}, which $\leftarrow$ Latin \emph{pensare}, \textquotedblleft{}to weigh, to ponder\textquotedblright{} (EtymDB 2.1) — the same Latin verb as \emph{pensar}.
\end{cousinweb}

\begin{grammarlens}[title={What you've built}]
One more word for this chapter.
\end{grammarlens}

\subsection*{Your turn}
\begin{itemize}
\raggedright
  \item \emph{Say it:} \emph{pesado}
  \item \emph{Say it:} \emph{pesado}, once more
  \item \emph{Say it:} say \emph{pesado}
  \item \emph{Recall:} say the Portuguese for salty, then the Portuguese for light, then say \emph{pesado} again
\end{itemize}

\begin{tcolorbox}[breakable,colback=teal!4,colframe=teal!35!black,title={Before you move on}]
What does pesado mean? (\textquotedblleft{}Heavy\textquotedblright{}.) Say it once more.
\end{tcolorbox}
\section[magro]{magro — thin}

\label{lesson:PT-C128-magro}



\begin{quote}
Before the new one: say the Portuguese for light, then the Portuguese for heavy.
\end{quote}

\subsection*{You'll want to know: magro}
\textbf{magro} — \textquotedblleft{}thin\textquotedblright{}.

\begin{cousinweb}
\textbf{magro} $\leftarrow$ Latin \emph{macer, macri}, \textquotedblleft{}lean\textquotedblright{} — inherited. EtymDB 2.1 traces it only to Old Portuguese \emph{magro}; it derives Spanish \emph{magro} from Latin \emph{macrum}.
\end{cousinweb}

\begin{grammarlens}[title={What you've built}]
One more word for this chapter.
\end{grammarlens}

\subsection*{Your turn}
\begin{itemize}
\raggedright
  \item \emph{Say it:} \emph{magro}
  \item \emph{Say it:} \emph{magro}, once more
  \item \emph{Say it:} say \emph{magro}
  \item \emph{Recall:} say the Portuguese for light, then the Portuguese for heavy, then say \emph{magro} again
\end{itemize}

\begin{tcolorbox}[breakable,colback=teal!4,colframe=teal!35!black,title={Before you move on}]
What does magro mean? (\textquotedblleft{}Thin\textquotedblright{}.) Say it once more.
\end{tcolorbox}
\section[contente]{contente — pleased, glad}

\label{lesson:PT-C128-contente}



\begin{quote}
Before the new one: say the Portuguese for heavy, then the Portuguese for thin.
\end{quote}

\subsection*{You'll want to know: contente}
\textbf{contente} — \textquotedblleft{}pleased, glad\textquotedblright{}.

\begin{cousinweb}
\textbf{contente} $\leftarrow$ Latin \emph{contentus}, \textquotedblleft{}held in, satisfied\textquotedblright{}, from \emph{continere}, \textquotedblleft{}to hold together\textquotedblright{} (Lewis \& Short). EtymDB 2.1 traces it only to Old Portuguese \emph{contente}.
\end{cousinweb}

\begin{grammarlens}[title={What you've built}]
One more word for this chapter.
\end{grammarlens}

\subsection*{Your turn}
\begin{itemize}
\raggedright
  \item \emph{Say it:} \emph{contente}
  \item \emph{Say it:} \emph{contente}, once more
  \item \emph{Say it:} say \emph{contente}
  \item \emph{Recall:} say the Portuguese for heavy, then the Portuguese for thin, then say \emph{contente} again
\end{itemize}

\begin{tcolorbox}[breakable,colback=teal!4,colframe=teal!35!black,title={Before you move on}]
What does contente mean? (\textquotedblleft{}Pleased, glad\textquotedblright{}.) Say it once more.
\end{tcolorbox}
\section[o fim de semana]{o fim de semana — the weekend}

\label{lesson:PT-C128-fim-de-semana}



\begin{quote}
Before the new one: say the Portuguese for thin, then the Portuguese for pleased.

The weekend is two days: say \emph{sábado} and \emph{domingo}.
\end{quote}

\subsection*{You'll want to know: o fim de semana}
\textbf{o fim de semana} — \textquotedblleft{}the weekend\textquotedblright{}.

\begin{cousinweb}
\textbf{fim de semana} is \textquotedblleft{}end of week\textquotedblright{}: \emph{fim} $\leftarrow$ Latin \emph{finis}, \textquotedblleft{}end\textquotedblright{} (Lewis \& Short), and \emph{semana} $\leftarrow$ Late Latin \emph{septimana} (EtymDB 2.1).
\end{cousinweb}

\begin{grammarlens}[title={What you've built}]
One more word for this chapter.
\end{grammarlens}

\subsection*{Your turn}
\begin{itemize}
\raggedright
  \item \emph{Say it:} \emph{o fim de semana}
  \item \emph{Say it:} \emph{o fim de semana}, once more
  \item \emph{Say it:} say \emph{o fim de semana}
  \item \emph{Recall:} say the Portuguese for thin, then the Portuguese for pleased, then say \emph{o fim de semana} again
\end{itemize}

\begin{tcolorbox}[breakable,colback=teal!4,colframe=teal!35!black,title={Before you move on}]
What does o fim de semana mean? (\textquotedblleft{}The weekend\textquotedblright{}.) Say it once more.
\end{tcolorbox}
\section[a data]{a data — a date}

\label{lesson:PT-C128-data}



\begin{quote}
Before the new one: say the Portuguese for pleased, then the Portuguese for the weekend.

Say the two weekend days again: \emph{sábado, domingo}.
\end{quote}

\subsection*{You'll want to know: a data}
\textbf{a data} — \textquotedblleft{}a date\textquotedblright{}.

\begin{cousinweb}
\textbf{data} $\leftarrow$ Late Latin \emph{data}, \textquotedblleft{}given\textquotedblright{} (EtymDB 2.1) — a learned borrowing from Late Latin. Letters began \emph{data}, \textquotedblleft{}given\textquotedblright{}, at a place on a day.
\end{cousinweb}

\begin{grammarlens}[title={What you've built}]
One more word for this chapter.
\end{grammarlens}

\subsection*{Your turn}
\begin{itemize}
\raggedright
  \item \emph{Say it:} \emph{a data}
  \item \emph{Say it:} \emph{a data}, once more
  \item \emph{Say it:} say \emph{a data}
  \item \emph{Recall:} say the Portuguese for pleased, then the Portuguese for the weekend, then say \emph{a data} again
\end{itemize}

\begin{tcolorbox}[breakable,colback=teal!4,colframe=teal!35!black,title={Before you move on}]
What does a data mean? (\textquotedblleft{}A date\textquotedblright{}.) Say it once more.
\end{tcolorbox}
